\section{Executive Summary}

\subsection{The problem in one paragraph}

An engineer at a refinery needs to read a scanned inspection report, cross-check
it against the governing standard operating procedure, work out which downstream
units are affected, and draft an approval note. Every artefact involved --- the
piping and instrumentation diagram, the vendor correspondence, the financials,
the unreleased design --- is confidential. Company policy forbids sending it to
a cloud assistant. So the engineer either does the work by hand, losing hours,
or pastes the material into a public tool anyway, creating exactly the exposure
the policy exists to prevent.

\subsection{What was built}

A complete, working system that runs on one consumer GPU with no network
connection. It comprises:

\begin{itemize}
\item a \term{model gateway} serving five open-weight models with automatic,
      VRAM-aware selection driven by a declarative manifest;
\item an \term{agent} that plans multi-step work, calls local tools, critiques
      its own output and iterates;
\item a \term{retrieval layer} combining dense vectors, sparse BM25 and a
      typed \term{knowledge graph}, with four retrieval strategies chosen by
      query shape;
\item a \term{multimodal ingestion pipeline} handling text PDFs, scanned pages,
      handwriting and engineering drawings;
\item a \term{deliverable generator} producing genuine \code{.docx},
      \code{.pptx} and \code{.xlsx} files, not chat transcripts;
\item a \term{sovereignty subsystem} providing kernel-level network isolation,
      a live egress monitor, a deliberate breach canary, and a hash-chained
      tamper-evident audit log.
\end{itemize}

\subsection{Status at the time of writing}

\begin{center}
\renewcommand{\arraystretch}{1.25}
\begin{tabularx}{\linewidth}{@{}l r X@{}}
\toprule
\textbf{Measure} & \textbf{Value} & \textbf{Note} \\
\midrule
Python source          & 3\,124 lines & 46 modules across 9 packages \\
Automated tests        & 61 passing   & full suite runs in under 1 second \\
Dependencies           & 99 packages  & all inside a project-local virtualenv \\
Ontology               & 23 node types & 31 relation types, 43 endpoint signatures \\
Open-weight models     & 5            & approximately 24\,GB of weights \\
Knowledge graph        & 499 nodes    & 391 typed relations from a 5-document corpus \\
Community summaries    & 37           & enables corpus-wide thematic queries \\
Verified rubric points & 5 of 5       & each demonstrated, not asserted \\
\bottomrule
\end{tabularx}
\end{center}

\subsection{The three claims that distinguish this work}

\begin{keybox}{1. Sovereignty is demonstrated, not claimed}
The problem statement is explicit that proof --- not assertion --- is what
counts. Generated code executes inside an unprivileged network namespace where
a connection attempt to \code{1.1.1.1:53} returns
\code{OSError: Network is unreachable}. This is a kernel guarantee, not a
policy setting, and it was verified experimentally.
\end{keybox}

\begin{keybox}{2. Graph RAG answers questions vector search structurally cannot}
``If we isolate P-101A, which downstream units are affected?'' requires
traversal. Similarity is not connectivity. A two-hop dependency shares no
vocabulary with the query, so an embedding model cannot retrieve it, whereas a
graph walk finds it deterministically. A P\&ID is \emph{already} a graph; the
system recovers that structure rather than flattening it into vectors.
\end{keybox}

\begin{keybox}{3. The domain ontology is configuration, not code}
The knowledge-graph schema lives in \code{config/ontology.yaml}. A
\code{general} pack covers any corpus; an \code{industrial} pack adds refinery
types. Packs merge, and a new domain is a YAML block rather than a code change.
This was validated by indexing a corpus of curricula vitae and a journal
article, which produced 499 nodes across Person, Concept, Technology,
Publication and Skill types.
\end{keybox}
